%! AP Statistics — Unit 5, Lesson 5.1 — Homework
\documentclass[10pt]{article}
\usepackage{apstats-article}
\usepackage{apstats-boxes}
\newcommand{\TallMath}[1]{$\displaystyle #1\rule[-1.4em]{0pt}{3.2em}$}

\begin{document}

\pageheader{Unit 5, Lesson 5.1}{Homework}
\namedateperiod

\vspace{0.15in}

\begin{enumerate}

  \item For each situation, identify the \textbf{statistic} and the \textbf{parameter}.
        Write the appropriate notation for each ($\bar{x}$, $\mu$, $\hat{p}$, or $p$).

        \begin{enumerate}[label=\alph*.]
          \item A researcher records the resting heart rate of a random sample of 80
                adults. The mean heart rate in the sample is 72 bpm. The researcher
                wants to estimate the mean heart rate of all adults.
                \begin{itemize}
                  \item Statistic: \blank{5cm} \quad Notation: \blank{1.8cm}
                  \item Parameter: \blank{5cm} \quad Notation: \blank{1.8cm}
                \end{itemize}
          \item In a survey of 200 randomly selected college students, 64 said they
                exercise at least 3 times per week. The survey is meant to estimate the
                proportion of all college students who do so.
                \begin{itemize}
                  \item Statistic: \blank{5cm} \quad Notation: \blank{1.8cm}
                  \item Parameter: \blank{5cm} \quad Notation: \blank{1.8cm}
                \end{itemize}
        \end{enumerate}

  \item A botanist measures the heights of 15 randomly selected seedlings from a
        greenhouse and finds a mean height of 4.3 cm. Another researcher measures 15
        different randomly selected seedlings from the same greenhouse and gets 4.7 cm.

        \begin{enumerate}[label=\alph*.]
          \item Are these two measurements the same thing? Explain. \writelines{2}
          \item Is either measurement necessarily wrong? Why or why not? \writelines{2}
          \item What do you call the variation between these two measurements?
                \blank{5cm}
        \end{enumerate}

  \item A simulation takes 500 random samples of size $n = 40$ from a population with
        mean $\mu = 100$ and records each sample mean. The 500 sample means are then
        displayed in a dot-plot.

        \begin{enumerate}[label=\alph*.]
          \item What is the name of the distribution shown in the dot-plot?
                \blank{7cm}
          \item Would you expect the center of the dot-plot to be near 100? Explain.
                \writelines{2}
          \item If the sample size were increased to $n = 200$, how would you expect
                the spread of the dot-plot to change? \writelines{2}
        \end{enumerate}

\end{enumerate}

\begin{extensionbox}
\textbf{Extension (optional).} Find a news article that reports two different survey
results about the same topic (e.g., two polls about political approval ratings, consumer
confidence, or vaccine attitudes). Write a short paragraph (3--4 sentences) explaining
whether the difference between the two results could be due to sampling variability or
might indicate a real difference in the populations surveyed.
\end{extensionbox}

\begin{reflectionbox}
\textbf{Preview:} Lesson 5.2 revisits the normal distribution, which will be the key
model for describing the shape of sampling distributions. Review your notes on
z-scores and normal probabilities from Unit 1.
\end{reflectionbox}

\end{document}
